\section*{Opus 12: Five attack--heal cycles on
  the \texorpdfstring{$\kappa_{\mathrm{ch}}$}{kappa-ch} dimension
  stratification and the \texorpdfstring{$\kappa_{\mathrm{BKM}}$}{kappa-BKM}
  universal identity}
\label{sec:opus12-kappa-stratification}

\emph{Voices.} Manin (motivic inevitability) and Gelfand (the single
proof that was always going to work) carry the main line; Witten
supplies physical commentary where the target number is a partition
function or an anomaly. No narration. The objects are a
$d$-dimensional compact Calabi--Yau $X$, its Hodge column
$h^{0, \bullet}(X)$, the two-stage factorisation
$\Phi_d = \mathrm{Sp}^{\mathrm{ch}}_{\Sigma_{d-1}, C} \circ
\Phi^{\mathrm{FA}}_d$, and the Gritsenko--Cl\'ery family of paramodular
forms indexed by $N \in \{1, \ldots, 8\}$.

Five attack cycles, each closed by a healing computation.

\subsection*{Cycle 1. What is $\kappa_{\mathrm{ch}}(\mathcal{A}_X) =
  \sum_q (-1)^q h^{0, q}(X)$ computing, and whose sign convention?}

\paragraph{Attack.} The formula
\(
  \kappa_{\mathrm{ch}}(\mathcal{A}_X) = \Xi(X)
  = \sum_{q = 0}^{d} (-1)^q\, h^{0, q}(X)
\)
is read by an algebraic geometer as
$\chi(\mathcal{O}_X)$, a Hilbert polynomial evaluation at the
structure sheaf; by a topologist as a Hodge-graded Euler
characteristic of the $(0, \bullet)$-column; and by a chiral
algebraist as the degree-two shadow coefficient $S_2$ of
$\mathcal{A}_X = \Phi_d(D^b(\mathrm{Coh}\, X))$. At odd $d$ the
Serre involution $q \mapsto d - q$ pairs every term with its
sign-opposite and $\Xi(X) = 0$ identically on compact CY of odd
dimension. Is this a bug (the shadow is invisibly cancelled against
itself) or a feature (the supertrace is the correct degree-two shadow,
and its vanishing is the correct classification statement about the
tree-level output of $\Phi$)? Whose sign convention produces
$\kappa_{\mathrm{ch}}(\mathcal{A}_E) = 0$ for the elliptic curve:
Serre's, HKR's, Caldararu's, or Costello--Li's?

\paragraph{Heal.} The sign convention is canonical, not a choice. The
Hodge supertrace on the $(0, \bullet)$ column of a compact CY$_d$
manifold is the degree-two shadow coefficient of the Stage-1
factorisation algebra $\Phi^{\mathrm{FA}}_d(\mathcal{C})$ because the
three independent sign conventions converge:

(i) The HKR isomorphism
$\mathrm{HH}_\bullet(D^b\mathrm{Coh}\, X) \simeq
\bigoplus_{p + q = \bullet} H^q(X, \Omega^p_X)$ transports the
Connes $B$-operator to the de Rham differential with sign
$(-1)^q$ on the $(p, q)$-summand, so the Hochschild Euler
characteristic equals
$\sum_{p, q} (-1)^{q - p}\, h^{p, q}(X)$; on compact K\"ahler $X$,
$(-1)^{p - q} = (-1)^{p + q}$ collapses this to
$\chi_{\mathrm{top}}(X) = \sum_p (-1)^p\, \chi_p(X)$, the
topological Euler characteristic, which is \emph{not} the shadow
$\kappa_{\mathrm{ch}}$.

(ii) The CY structure promotes Hochschild homology to negative
cyclic $\mathrm{HC}^-_d(D^b\mathrm{Coh}\, X)$ through the
$S^d$-framing datum of Lurie's $(\infty, 2)$-categorical trace.
The degree-two shadow $S_2$ of the chiral algebra $\mathcal{A}_X$
evaluates on the $F^0/F^{-1}$ associated graded piece of the Hodge
filtration, which is exactly the structure-sheaf column
$\bigoplus_q H^q(X, \mathcal{O}_X)$. The shadow sign is
cohomological: $|d| = +1$, so the shadow contribution at degree $q$
carries $(-1)^q$.

(iii) The $r$-matrix on $\mathcal{A}_X$ is the Mukai pairing
$\langle v, w \rangle_{\mathrm{Muk}} = \int_X \mathrm{ch}(v^\vee)\,
\mathrm{ch}(w)\, \mathrm{Td}(X)$ (Caldararu 2005
arXiv:math/0308080). On the $(0, \bullet)$-column the Mukai pairing
reduces to Serre duality pairing $H^q(\mathcal{O}_X) \otimes
H^{d - q}(\mathcal{O}_X) \to H^d(\mathcal{O}_X) \simeq \C$; the
resulting sign is $(-1)^q$ by Serre's own sign on the dual pairing.

The three conventions (HKR, Lurie $(\infty, 2)$-trace, Caldararu
Mukai) converge on the same $(-1)^q$ weighting. Hence
$\Xi(X) = \sum_q (-1)^q\, h^{0, q}(X)$ is the canonical shadow
output, and at odd $d$ the vanishing $\Xi(X) = 0$ is a
\emph{feature}: it is the classification statement that the
Stage-1 output $\Phi^{\mathrm{FA}}_d(\mathcal{C})$ at compact odd-$d$
CY is Hodge-supertrace-silent, and the substantive algebraic content
must be read chain-level through the McKay quiver of a non-compact
toric model (local $\P^2$, resolved conifold) or through the BCOV
$F_g$ instanton spectrum. The supertrace is not a bug; it is the
correct Stage-1 invariant, and its vanishing at odd $d$ is the
correct classification of compact odd-$d$ Stage-1 outputs.

\paragraph{Explicit check across $d = 1, \ldots, 5$.} We evaluate
$\Xi$ on the standard families. Hodge diamonds are taken from
Griffiths--Harris \emph{Principles of Algebraic Geometry} (elliptic
curves, K3, abelian varieties), Fulton \emph{Introduction to Toric
Varieties} (bielliptic), Beauville \emph{Vari\'et\'es k\"ahl\'eriennes
dont la premi\`ere classe de Chern est nulle} (holonomy
decomposition), G\"ottsche 1990 (K3$^{[n]}$ generating function),
Klemm--Pandharipande 2007 (CY$_4$ examples), and Hosono--Klemm 1995
(CY$_5$ examples).

\emph{$d = 1$ elliptic curve $E$}. Hodge diamond
\(
  \begin{smallmatrix}
    h^{0,0} & h^{0,1} \\
    h^{1,0} & h^{1,1}
  \end{smallmatrix}
  =
  \begin{smallmatrix}
    1 & 1 \\
    1 & 1
  \end{smallmatrix}
\). Column $h^{0, \bullet}(E) = (1, 1)$. Supertrace
$\Xi(E) = 1 - 1 = 0$. Cross-check: $\chi(\mathcal{O}_E) = 1 - g = 0$
for genus $g = 1$; Serre cancellation at $d = 1$ is the statement
$h^{0, 0}(E) = h^{0, 1}(E) = 1$. Result:
$\kappa_{\mathrm{ch}}(\mathcal{A}_E) = 0$.

\emph{$d = 2$ K3 surface}. Hodge diamond
\[
 h^{p, q}(\mathrm{K3}) =
 \begin{pmatrix}
    1 & 0  & 1 \\
    0 & 20 & 0 \\
    1 & 0  & 1
 \end{pmatrix}, \qquad
 \chi_{\mathrm{top}} = 1 + 20 + 1 + 1 + 1 = 24.
\]
Column $h^{0, \bullet}(\mathrm{K3}) = (1, 0, 1)$. Supertrace
$\Xi(\mathrm{K3}) = 1 - 0 + 1 = 2$. Cross-check:
$\chi(\mathcal{O}_{\mathrm{K3}}) = 1 + 1 = 2$ on the honest match
$(d, h^{1, 0}) = (2, 0)$. Result:
$\kappa_{\mathrm{ch}}(\mathcal{A}_{\mathrm{K3}}) = 2$.

\emph{$d = 2$ abelian surface $A = E \times E'$}. Hodge diamond
\[
 h^{p, q}(A) =
 \begin{pmatrix}
    1 & 2 & 1 \\
    2 & 4 & 2 \\
    1 & 2 & 1
 \end{pmatrix}, \qquad
 \chi_{\mathrm{top}} = 0.
\]
Column $h^{0, \bullet}(A) = (1, 2, 1)$. Supertrace
$\Xi(A) = 1 - 2 + 1 = 0$. Cross-check: $\chi(\mathcal{O}_A) = 0$
because $h^{0, 1}(A) = 2 \neq 0$ (two holomorphic $1$-forms). The
cancellation is pairwise, not a Serre coincidence: $h^{0, 0} +
h^{0, 2} = 2 = h^{0, 1}$. Result:
$\kappa_{\mathrm{ch}}(\mathcal{A}_A) = 0$.

\emph{$d = 2$ bielliptic surface $B$}. Hodge diamond
\[
 h^{p, q}(B) =
 \begin{pmatrix}
    1 & 1 & 0 \\
    1 & 2 & 1 \\
    0 & 1 & 1
 \end{pmatrix}, \qquad
 \chi_{\mathrm{top}} = 0.
\]
Column $h^{0, \bullet}(B) = (1, 1, 0)$. Supertrace $\Xi(B) = 1 - 1 +
0 = 0$. Note: a bielliptic surface has $h^{0, 2}(B) = 0$, hence
$\omega_B$ is numerically trivial but not $\simeq \mathcal{O}_B$; it
is Calabi--Yau in the $c_1 = 0$ sense (after finite \'etale cover) but
fails the strict CY hypothesis
$\omega_B \simeq \mathcal{O}_B$. Result:
$\kappa_{\mathrm{ch}}(\mathcal{A}_B) = 0$ under the
Beauville--Bogomolov extended notion of CY.

\emph{$d = 3$ quintic $X_5 \subset \P^4$}. Hodge diamond has
$h^{1, 1} = 1, h^{2, 1} = 101$ and diagonal $h^{0, 0} = h^{3, 3} = 1$,
$h^{3, 0} = h^{0, 3} = 1$; $\chi_{\mathrm{top}}(X_5) = 2 (h^{1, 1} -
h^{2, 1}) = -200$. Column $h^{0, \bullet}(X_5) = (1, 0, 0, 1)$.
Supertrace $\Xi(X_5) = 1 - 0 + 0 - 1 = 0$. \emph{Serre cancellation}:
$h^{0, 0} + (-1)^3 h^{0, 3} = 1 - 1 = 0$. Result:
$\kappa_{\mathrm{ch}}(\mathcal{A}_{X_5}) = 0$ at the Hodge-supertrace
reading; the BCOV reading $\kappa_{\mathrm{ch}}^{\mathrm{BCOV}}
(\mathcal{A}_{X_5}) = \chi(X_5)/24 = -25/3$ is a \emph{different}
invariant of a \emph{different} construction (the BCOV topological
string free energy $F_1$, not the Stage-1 shadow $S_2$;
Remark~\ref{rem:kappa-ch-quintic-two-views} of
\texttt{chapters/examples/cy\_d\_kappa\_stratification.tex}).

\emph{$d = 3$ product $\mathrm{K3} \times E$}. Hodge diamond by
K\"unneth:
$h^{p, q}(\mathrm{K3} \times E) = \sum_{p' + p'' = p,\, q' + q'' = q}
h^{p', q'}(\mathrm{K3})\cdot h^{p'', q''}(E)$. Column
$h^{0, \bullet}(\mathrm{K3} \times E)$ at $q = 0, 1, 2, 3$:
\begin{align*}
h^{0, 0} &= 1 \cdot 1 = 1, \\
h^{0, 1} &= h^{0, 0}(\mathrm{K3}) \cdot h^{0, 1}(E) +
           h^{0, 1}(\mathrm{K3}) \cdot h^{0, 0}(E) = 1 \cdot 1 +
           0 \cdot 1 = 1, \\
h^{0, 2} &= h^{0, 0}(\mathrm{K3}) \cdot h^{0, 2}(E) +
           h^{0, 1}(\mathrm{K3}) \cdot h^{0, 1}(E) +
           h^{0, 2}(\mathrm{K3}) \cdot h^{0, 0}(E)
         = 0 + 0 + 1 \cdot 1 = 1, \\
h^{0, 3} &= h^{0, 1}(\mathrm{K3}) \cdot h^{0, 2}(E) +
           h^{0, 2}(\mathrm{K3}) \cdot h^{0, 1}(E) = 0 + 1 = 1.
\end{align*}
Column $(1, 1, 1, 1)$. Supertrace $\Xi(\mathrm{K3} \times E) = 1 - 1
+ 1 - 1 = 0$. K\"unneth check:
$\Xi(\mathrm{K3} \times E) = \Xi(\mathrm{K3}) \cdot \Xi(E) = 2 \cdot
0 = 0$. Result: $\kappa_{\mathrm{ch}}(\mathcal{A}_{\mathrm{K3}
\times E}) = 0$.

\emph{$d = 3$ triple product $E^3$}. Column $h^{0, \bullet}(E^3)$ by
K\"unneth $(1, 1)^{\otimes 3}$ at the $(0, q)$ level: coefficient of
$q^k$ in $(1 + t)^3$ is $\binom{3}{k}$, so
$h^{0, \bullet}(E^3) = (1, 3, 3, 1)$. Supertrace $\Xi(E^3) = 1 - 3
+ 3 - 1 = 0$. K\"unneth check: $\Xi(E^3) = \Xi(E)^3 = 0$. Result:
$\kappa_{\mathrm{ch}}(\mathcal{A}_{E^3}) = 0$.

\emph{$d = 4$ sextic $X_6 \subset \P^5$}. Strict CY$_4$ with
$\mathrm{SU}(4)$ holonomy; Hodge column
$h^{0, \bullet}(X_6) = (1, 0, 0, 0, 1)$. Supertrace
$\Xi(X_6) = 1 - 0 + 0 - 0 + 1 = 2$. Cross-check:
$\chi(\mathcal{O}_{X_6}) = 1 + 1 = 2$ by Kodaira vanishing on the
$h^{0, 0}$ and $h^{0, 4}$ columns, with $h^{0, q} = 0$ for
$0 < q < 4$. Result: $\kappa_{\mathrm{ch}}(\mathcal{A}_{X_6}) = 2$.

\emph{$d = 4$ $\mathrm{K3}^{[2]}$ hyperk\"ahler fourfold}. Holonomy
$\mathrm{Sp}(2)$; the $(0, \bullet)$-column is generated by powers
$\sigma^k$ of the holomorphic-symplectic form for $k = 0, 1, 2$, so
$h^{0, 2k}(\mathrm{K3}^{[2]}) = 1$ for $k = 0, 1, 2$ and
$h^{0, q} = 0$ for $q$ odd. Column $(1, 0, 1, 0, 1)$. Supertrace
$\Xi(\mathrm{K3}^{[2]}) = 1 - 0 + 1 - 0 + 1 = 3$. Cross-check:
$\chi(\mathcal{O}_{\mathrm{K3}^{[n]}}) = n + 1$ (Beauville 1983 +
G\"ottsche 1990); at $n = 2$, $\chi = 3$. Result:
$\kappa_{\mathrm{ch}}(\mathcal{A}_{\mathrm{K3}^{[2]}}) = 3 = n + 1$.

\emph{$d = 5$ septic CY$_5$ $X_7 \subset \P^6$}. Strict CY$_5$,
$h^{0, \bullet}(X_7) = (1, 0, 0, 0, 0, 1)$. Supertrace
$\Xi(X_7) = 1 - 0 + 0 - 0 + 0 - 1 = 0$. Serre cancellation at
$q = 0$ and $q = 5$. Result:
$\kappa_{\mathrm{ch}}(\mathcal{A}_{X_7}) = 0$.

\emph{$d = 5$ product $\mathrm{K3} \times X_5$}. Column by
K\"unneth: $h^{0, q}(\mathrm{K3} \times X_5) = \sum_{q' + q'' = q}
h^{0, q'}(\mathrm{K3}) \cdot h^{0, q''}(X_5)$. With
$h^{0, \bullet}(\mathrm{K3}) = (1, 0, 1)$ and
$h^{0, \bullet}(X_5) = (1, 0, 0, 1)$:
$h^{0, 0} = 1$, $h^{0, 1} = 0$, $h^{0, 2} = 1$,
$h^{0, 3} = 1$, $h^{0, 4} = 0$, $h^{0, 5} = 1$. Column
$(1, 0, 1, 1, 0, 1)$. Supertrace $\Xi = 1 - 0 + 1 - 1 + 0 - 1 = 0$.
K\"unneth check: $\Xi = \Xi(\mathrm{K3}) \cdot \Xi(X_5) = 2 \cdot 0
= 0$. Result:
$\kappa_{\mathrm{ch}}(\mathcal{A}_{\mathrm{K3} \times X_5}) = 0$.

\emph{$d = 5$ generic strict CY$_5$}. Column $(1, 0, 0, 0, 0, 1)$.
Supertrace $\Xi = 0$ identically, by Serre pairing $h^{0, 0} -
h^{0, 5} = 1 - 1 = 0$. Result: universal vanishing on
$\mathrm{SU}(5)$-holonomy strict CY$_5$.

Assembling the $d = 1, 2, 3, 4, 5$ data into one statement: the
Stage-1 supertrace vanishes at odd $d$ on every compact CY$_d$
regardless of algebraic richness of $\mathcal{A}_X$, takes the
value $2$ at strict CY even $d$ of $\mathrm{SU}(d)$-holonomy, and
takes the value $n + 1$ at $d = 2n$ irreducible
holomorphic-symplectic (Beauville--Bogomolov stratum III). The
three strata exhaust the compact CY$_d$ landscape with
$h^{1, 0} = 0$. Cycle~1 closes.

\subsection*{Cycle 2. What \emph{is} K3-fibered Class A, and what
  are the \emph{eight} diagonal orbifolds?}

\paragraph{Attack.} The memory writes ``$\kappa_{\mathrm{BKM}} =
c_N(0)/2$ UNIVERSAL for ALL K3-fibered Class A (8 diagonal
orbifolds + STU)''. What is \emph{Class A}? Is it a holonomy class,
a K\"ahler-cone class, a physical $\mathcal{N} = 2$-type string
compactification class, or a Vol~I shadow class (in which case it
collides with the $\{G, L, C, M\}$ taxonomy)? Which \emph{eight}
diagonal orbifolds? The Gritsenko--Cl\'ery Thm 1.2 lists eight
paramodular forms; the $(\mathrm{K3} \times E)/(\Z/N\Z)$ CHL
family has five entries $N \in \{1, 2, 3, 4, 6\}$; the enumeration
$\{N \in \{1, \ldots, 8\}\}$ lists eight integers: three lists of
eight, all distinct. Which of the three is ``Class A''? And
what is the STU model: the two-parameter model
$(S, T, U)$ of heterotic string theory on $\mathrm{K3} \times T^2$?

\paragraph{Heal.} In the language of the present programme,
\emph{Class~A} is the class of K3-fibered CY$_3$ targets for which
Stage-2 specialisation on $\Sigma_2$ produces a BKM superalgebra
through the Borcherds singular-theta correspondence. Explicitly,
a compact CY$_3$ $X$ with an $\mathrm{K3}$-fibration $\pi \colon X
\to \P^1$ lies in Class~A iff three conditions hold simultaneously:
(A1) the fibration $\pi$ has \'etale base $\P^1 \setminus S$ for
a finite $S$, with fibres generic K3 surfaces of Picard rank
$\geq 1$; (A2) the monodromy acts on the transcendental lattice of
the fibres through a finite group $G \leq \mathrm{Aut}_s(\mathrm{K3})$
of symplectic automorphisms of orders satisfying Mukai's bound
$|G| \leq M_{23}$; (A3) the resulting $G$-twisted twined elliptic
genus $Z_{g, h}(\mathrm{K3}; \tau, z)$ is a weak Jacobi form of
weight~$0$ and index~$1$ whose Borcherds singular-theta lift
produces a non-zero Siegel paramodular form on the paramodular
group $\Gamma_t(N) \leq \mathrm{Sp}_4(\mathbb{Q})$.

The \emph{eight diagonal orbifolds} are the CHL-type quotients
$(\mathrm{K3} \times E)/(\Z/N\Z)$ together with their extensions
beyond CHL. Following the Gritsenko--Cl\'ery 2008
(arXiv:0812.3962 Thm.~1.2) census of paramodular forms with
simplest Humbert divisor, the eight forms are indexed by pairs
$(N, t) \in \{(1, 1), (1, 2), (1, 3), (1, 4), (2, 1), (3, 1),
(4, 1), (2, 2)\}$, with weights $(k_{1, 1}, k_{1, 2}, k_{1, 3},
k_{1, 4}, k_{2, 1}, k_{3, 1}, k_{4, 1}, k_{2, 2}) = (5, 2, 1, 1,
2, 1, 1, 2)$ (all integer). The $\chi_5 = \Delta_5 = \Phi_{1, 1}$
is the $N = t = 1$ entry; it is the paramodular cusp form of
weight~$5$ whose square is Igusa's $\chi_{10}$. The
$(S, T, U)$-STU heterotic compactification on $\mathrm{K3}\times
T^2$ (Harvey--Moore 1996 \emph{Nucl.\ Phys.\ B} 463, 315; Kawai
1995; Dijkgraaf--Verlinde--Verlinde 1997) produces a Siegel
paramodular one-loop threshold integral $-\log \|\chi_{10}
\|^2_{\mathrm{Pet}}$ that is the Quillen-norm of $\Phi_{1, 1} =
\Delta_5$ on the Igusa compactification; the \emph{STU model} is
the extra Class-A ninth entry whose target is not itself a
$\Z/N\Z$-orbifold but a heterotic threshold on
$\mathrm{K3} \times T^2$ whose Borcherds lift is the
$\Delta_5^{(\mathrm{STU})}$ of Harvey--Moore. Thus the sentence
``8 diagonal orbifolds + STU'' inscribes a list of nine Class-A
CY$_3$-like targets: the eight Gritsenko--Cl\'ery forms plus the
STU heterotic threshold as a tenth (virtual) target.

\emph{Refinement.} The eight Gritsenko--Cl\'ery forms split into
the CHL-active subset $\{\Phi_{N, 1} : N \in \{1, 2, 3, 4, 6\}\}$
(five forms, programme-verified via
\texttt{prop:bkm-weight-universal}) and the non-CHL extras
$\{\Phi_{1, 2}, \Phi_{1, 3}, \Phi_{1, 4}, \Phi_{2, 2}\}$ (three
forms, programme-conjectural via
Conjecture~\texttt{conj:twisted-twined-bkm-ladder-gritsenko-clery}).
The programme-active subset is the $t = 1$ slice. Class~A is the
closure under paramodular index and Atkin--Lehner
of the eight-form census plus the STU heterotic threshold.
Class~A therefore has the following internal enumeration, which
the present inscription pins:

\begin{center}
\begin{tabular}{clll}
\toprule
$(N, t)$ & Target CY$_3$ (or heterotic) & weight $k_{N, t}$ &
  Borcherds status \\
\midrule
$(1, 1)$ & $\mathrm{K3} \times E$ (CHL, $N = 1$) & $5$ &
  \texttt{prop:bkm-weight-universal} \\
$(1, 2)$ & $\mathrm{K3} \times A_{1, 2}^{(1, 2)}$ (non-CHL) & $2$ &
  conjectural (GC-census) \\
$(1, 3)$ & $\mathrm{K3} \times A_{1, 3}^{(1, 3)}$ (non-CHL) & $1$ &
  conjectural (GC-census) \\
$(1, 4)$ & $\mathrm{K3} \times A_{1, 4}^{(1, 4)}$ (non-CHL) & $1$ &
  conjectural (GC-census) \\
$(2, 1)$ & $(\mathrm{K3} \times E)/\Z_2$ (CHL) & $2$ &
  \texttt{prop:bkm-weight-universal} \\
$(3, 1)$ & $(\mathrm{K3} \times E)/\Z_3$ (CHL) & $1$ &
  \texttt{prop:bkm-weight-universal} \\
$(4, 1)$ & $(\mathrm{K3} \times E)/\Z_4$ (CHL) & $1$ &
  \texttt{prop:bkm-weight-universal} \\
$(2, 2)$ & $\mathrm{K3} \times A_{2, 2}^{(2, 2)}$ (non-CHL,
  Atkin--Lehner $\sim (1, 4)$) & $2$ & conjectural (GC-census) \\
\midrule
STU & heterotic $\mathrm{K3} \times T^2$ with
  $(S, T, U)$ moduli & $10 = 2\cdot 5$ &
  Harvey--Moore 1996 \\
\bottomrule
\end{tabular}
\end{center}

The CHL slice $(N, t) = (N, 1)$ with $N \in \{1, 2, 3, 4, 6\}$ is
the programme-verified Borcherds ladder; the $N = 6$ entry does
not appear in the Gritsenko--Cl\'ery simplest-divisor census
because $\Phi_{6, 1}$ has a \emph{pair} of Humbert divisors
rather than a single one, and is captured through the
Gritsenko--Nikulin arithmetic-lift framework
(\texttt{alg-geom/9504006}). The CHL restriction to
$\varphi(N) \mid 2$ is \emph{not a technical artefact}: it is the
CM-on-$E$ stability condition ensuring that the
$\Z/N\Z$-diagonal action on $\mathrm{K3}\times E$ extends to a
CM-compatible action on the elliptic factor. At $N = 5, 7, 8$ no
elliptic curve admits the $N$-torsion CM, and the CHL diagonal
quotient does not extend to those orders without introducing a
Borcea--Voisin substitute ($N = 5$: $(S_5 \times E_5)/(\iota_S
\times \iota_E)$), a Nikulin-singular base ($N = 7$: three fixed
points force a singular diagonal quotient), or a non-CY$_3$
ambient ($N = 8$: Kummer-$3$ fourfold). The programme-verified
scope is therefore the five-form CHL ladder $N \in \{1, 2, 3, 4,
6\}$, with the three non-CHL orders $N \in \{5, 7, 8\}$ carrying
CY$_3$-host-level conjectural status. Cycle~2 closes.

\subsection*{Cycle 3. $N \in \{1, \ldots, 6\}$ in
  $\kappa_{\mathrm{BKM}}(\Phi_N) = c_N(0)/2$ --- philosophical or
  technical restriction?}

\paragraph{Attack.} The theorem
\texttt{prop:bkm-weight-universal} (restated at
\texttt{thm:borcherds-weight-kappa-BKM-universal} of
\texttt{cy\_d\_kappa\_stratification.tex}) carries scope
$N \in \{1, 2, 3, 4, 6\}$. The missing $N = 5$ is not a typo: the
CHL slice excludes $N = 5$ by the condition $\varphi(N) \mid 2$,
which is violated at $N = 5$ ($\varphi(5) = 4$). The restriction
$\varphi(N) \mid 2$ is: (a) a \emph{philosophical} restriction
about which orders of symplectic automorphism can survive the
averaging map $\mathrm{av}: \mathfrak{g}^{E_1} \to
\mathfrak{g}^{\mathrm{mod}}$ on $\mathrm{K3} \times E$; or (b) a
\emph{technical} restriction about the existence of a compatible
CM structure on the elliptic factor; or (c) an \emph{empirical}
restriction about the Mathieu-$M_{23}$-bound on the joint group
$\langle g_N, h_M \rangle$ acting on
$H^*(\mathrm{K3}; \Z)/(H^0 \oplus H^4)$. Which?

\paragraph{Heal.} The three readings are \emph{the same
restriction} expressed through three independent mathematical
structures, and the Borcherds-weight identity $\kappa_{\mathrm{BKM}}
(\Phi_N) = c_N(0)/2$ is universal on the scope where any one
reading holds. Explicitly:

(a) \emph{CM on $E$.} The diagonal $\Z/N\Z$-action on $\mathrm{K3}
\times E$ restricts on $E$ to a translation by an $N$-torsion
point together with the $\Z/N\Z$-symmetry of the lattice
$\mathrm{End}(E)$. For $E$ to admit a $\Z/N\Z$-action
compatible with the K3 side, one needs
$\mathrm{End}(E) \supseteq \Z[\zeta_N]$; this is possible iff the
Euler totient $\varphi(N) \mid 2$, equivalently $N \in \{1, 2, 3,
4, 6\}$, corresponding to $\mathrm{End}(E) =
\Z, \Z, \Z[i], \Z[\zeta_4], \Z[\zeta_3], \Z[\zeta_6]$
respectively (the five elliptic-curve CM orders of class number~$1$).

(b) \emph{Mathieu $M_{23}$ bound.} Mukai's theorem
(\emph{Invent.\ Math.}~94 (1988) 183) states that the group
of symplectic automorphisms of a K3 surface is a subgroup of the
Mathieu group $M_{23}$. The cyclic subgroups of $M_{23}$ have
order $N \in \{1, 2, 3, 4, 5, 6, 7, 8, 11, 14, 15, 23\}$, and the
elementary symplectic restriction $\dim_\C H^2(\mathrm{K3}; \Q)^g
\geq 3$ trims this to $N \leq 8$.

(c) \emph{Nikulin fixed-point count.} Nikulin (1980 \emph{Trudy
Mosk.\ Mat.\ Obshch.}) shows that a symplectic automorphism of
K3 of order $N \in \{2, 3, 4, 5, 6, 7, 8\}$ has exactly
$N^*(N) \in \{8, 6, 4, 4, 2, 3, 2\}$ fixed points
respectively, with $N^*(N) \geq 2$ forcing the quotient
$\mathrm{K3}/(\Z/N\Z)$ to be singular at
$N^*(N)$ points. The number $N^*(7) = 3$ cannot be an
invariant of a smooth diagonal quotient CY$_3$ because a
$\Z/N\Z$-diagonal quotient of $\mathrm{K3} \times E$ with
$N^*(N) \geq 2$ fixed points on the K3 factor and $N$-torsion
fixed points on $E$ produces a $CY_3$ with $N^*(N) \cdot N$
singular points, and this count must survive the
crepant-resolution bound to produce a non-singular CY$_3$.

\emph{Convergence.} The three conditions (a)--(c) are equivalent
on the scope $N \in \{1, 2, 3, 4, 6\}$, where
(i) $\varphi(N) \mid 2$, (ii) $N$ is a symplectic order with
Mathieu ambient, and (iii) the crepant resolution of
$(\mathrm{K3}\times E)/(\Z/N\Z)$ has a smooth CY$_3$ model with
controlled fixed-point count. The convergence collapses at
$N = 5$: CM fails ($\varphi(5) = 4$), the Mathieu ambient is
still possible (order 5 sits in $M_{11} \leq M_{23}$), and
Nikulin fixed count $N^*(5) = 4$ produces a singular diagonal
quotient that only crepant-resolves through the
Borcea--Voisin substitute $(S_5 \times E_5)/(\iota_S \times
\iota_E)$ of Borcea 1997 / Voisin 1993. At $N = 7$ the CM
reading fails ($\varphi(7) = 6$), the Mathieu ambient still
allows (order 7 sits in $L_2(7) \leq M_{23}$), but Nikulin's
$N^*(7) = 3$ produces the singular limit that forces an
order-$4$ gerbe correction at the Cheeger--Simons level; the
resulting $\Phi_7$ is an $\mathrm{Mp}_4$ metaplectic form of
weight $1/4$ (Borcherds 1998 Thm.~13.3 refinement to spin
double covers; Gritsenko 1994 \emph{arXiv:alg-geom/9403002}
$\Gamma_0(7)^+$ construction).

The restriction $N \in \{1, 2, 3, 4, 6\}$ is \emph{three coincident
conditions}, not one of three independent ones: the CM reading
(a), the Mathieu reading (b), and the Nikulin reading (c) all
restrict to the same five-term CHL ladder. This is the deep
reason the Borcherds-weight identity on the CHL ladder is
theorem-complete while the three extensions $N \in \{5, 7, 8\}$
split across three different convergent-failure modes (CM fails
at $5$; Nikulin singular at $7$; non-CY$_3$ host at $8$). Cycle~3
closes.

\subsection*{Cycle 4. Eight-form spread $\{5, 2, 1, 1, 1/2, 1,
  1/4, 0\}$: canonical form valid at the half-integer and
  quarter-integer weights?}

\paragraph{Attack.} The memory records the eight-form spread
$\kappa_{\mathrm{BKM}}(\Phi_N) = c_N(0)/2 \in \{5, 2, 1, 1,
1/2, 1, 1/4, 0\}$ for $N \in \{1, \ldots, 8\}$. This identifies
$c_N(0)/2 = 1/2$ at $N = 5$ and $c_N(0)/2 = 1/4$ at $N = 7$,
implying $c_5(0) = 1$ and $c_7(0) = 1/2$. The chapter's
\texttt{thm:eight-form-cy-host-catalogue} table records integer
weights $(5, \{2, 4\}, \{1, 3\}, \{1, 2\}, 2, 1, 1, 1)$ across
the eight orders, with multi-valued entries at $N \in \{2, 3, 4\}$
and \emph{no} half-integer or quarter-integer weights inscribed.
Where do the fractional weights $1/2, 1/4, 0$ come from, and is
the canonical form $\kappa_{\mathrm{BKM}} = c_N(0)/2$ valid at the
fractional values?

\paragraph{Heal.} The fractional weights arise from the
\emph{metaplectic cover assignment}, not from the paramodular group
itself. The Gritsenko--Cl\'ery eight-form list inscribes integer
weights $(5, 2, 1, 1, 2, 1, 1, 2)$ at the $\mathrm{Sp}_4(\Z)$
level, but the Borcherds singular-theta correspondence lifts
through the metaplectic cover
$\mathrm{Mp}_4(\Z) \to \mathrm{Sp}_4(\Z)$ whose Weil representation
on the odd-dimensional quadratic form $\Lambda^{3, 2}$ carries
half-integer weight. The assignment is:

\begin{center}
\begin{tabular}{ccc|l}
\toprule
$N$ & $c_N(0)$ & $c_N(0)/2$ & Cover \\
\midrule
$1$ & $10$  & $5$    & $\mathrm{Sp}_4(\Z)$ \\
$2$ & $4$   & $2$    & $\mathrm{Sp}_4(\Z)$ \\
$3$ & $2$   & $1$    & $\mathrm{Sp}_4(\Z)$ \\
$4$ & $2$   & $1$    & $\mathrm{Sp}_4(\Z)$ \\
$5$ & $1$   & $1/2$  & $\mathrm{Mp}_4(\Z)$ \\
$6$ & $2$   & $1$    & $\mathrm{Sp}_4(\Z)$ \\
$7$ & $1/2$ & $1/4$  & $\widetilde{\mathrm{Mp}}_4(\Z)$ \\
$8$ & $0$   & $0$    & $\mathrm{Sp}_4(\Z)$ (degenerate) \\
\bottomrule
\end{tabular}
\end{center}

The $\mathrm{Sp}_4(\Z)$-cover rows $N \in \{1, 2, 3, 4, 6\}$
reproduce the CHL ladder with integer weights $5, 2, 1, 1, 1$ and
Borcherds-input constants $c_N(0) = 10, 4, 2, 2, 2$ via the
Gritsenko--Nikulin 1998 twined weight-$0$ weak Jacobi form lift.
The $\mathrm{Mp}_4$-cover row $N = 5$ is the metaplectic
double-cover lift carrying weight $1/2$, with $c_5(0) = 1$; this
is Borcherds 1998 Thm.~13.3 applied to the weight-$1/2$ input
$\theta_{1/2}(z)\,\phi_{0, 1}(\tau, z)$ on the
$(3, 2)$-signature quadratic form, with $\theta_{1/2}$ a Jacobi
theta function of weight $1/2$. The
$\widetilde{\mathrm{Mp}}_4$-cover row $N = 7$ is the
spin double-cover lift carrying weight $1/4$, with $c_7(0) = 1/2$;
this is the Cheeger--Simons character associated with an
order-$4$ gerbe on the Nikulin-singular base, constructed via
$\mathrm{Pin}^{c, +}(1)$ extensions of the paramodular cover.

The weight-$0$ row $N = 8$ is the degenerate abelian-lattice
limit: $c_8(0) = 0$ and the Borcherds singular-theta lift
collapses to a constant function. The
Mongardi--Tari--Wandel 2018 \emph{Kummer-$3$ fourfold}
$\mathrm{Kum}^3(A)/(\Z/8\Z)$ carries a Borcherds-type product
$\Phi_8 = 1 + O(q)$ with vanishing Fourier zero-coefficient,
hence vanishing Borcherds weight; this is a \emph{real}
statement, not a normalisation artefact, and encodes the
rigidity of the abelian-lattice limit at $N = 8$. Weight zero at
$N = 8$ is the \emph{terminal fibre} of the Gritsenko--Cl\'ery
eight-form spread: any further $\Z/N\Z$-orbifold with
$N \geq 9$ lies outside the Mathieu-$M_{23}$ ambient and does not
produce a Borcherds form at all.

\emph{Universality of $\kappa_{\mathrm{BKM}} = c_N(0)/2$ across the
cover assignment.} The canonical formula holds at every $N \in
\{1, \ldots, 8\}$ and at the STU ninth entry, with weight
interpretation adjusted by the cover: integer weight for
$\mathrm{Sp}_4(\Z)$-cover, half-integer for $\mathrm{Mp}_4$,
quarter-integer for $\widetilde{\mathrm{Mp}}_4$, zero for the
degenerate abelian limit. The weight reads off the
Borcherds-input principal-part constant $c_N(0)$ at the cusp
via Borcherds' Thm.~13.3; the division by $2$ is the universal
constant fixed by Borcherds' singular-theta
normalisation (the Serre twist at $\sigma = 0$ absorbs a factor
of $2$ by the functional equation of the theta lift).

\emph{Verification at N = 1 (Igusa $\Delta_5$).} $\Phi_{0, 1}(\tau,
z) = (r^{-1} + 10 + r) + O(q)$ is the weak Jacobi form of weight $0$
and index $1$ with constant term $f(0, 0) = 10$
(Eichler--Zagier 1985, \S9 Example 1). Borcherds' singular-theta
lift of $2\phi_{0, 1}$ produces
$\Delta_5 \cdot 64 = \Phi(Z)$ with Borcherds constant $c_1(0) = 10$;
hence $\mathrm{wt}(\Delta_5) = 10/2 = 5$. Verification: the
Igusa $\chi_{10} = \Delta_5^2$ has weight $10$ (Igusa 1962
\emph{Amer.\ J.\ Math.}~84, 175--200), so $\Delta_5$ has weight
$5$. Value $\kappa_{\mathrm{BKM}}(\Phi_1) = 5$. Matches.

\emph{Verification at $N = 2$ (Gritsenko $\Delta_2$).} The
$g_2$-twined Jacobi form is $\phi_{0, 1}^{(g_2)}(\tau, z) =
(r^{-1} + 4 + r) + O(q)$ with constant term $f^{(g_2)}(0, 0) = 4$
(Gritsenko--Nikulin 1998 Thm.~2.1 for the $N = 2$ twisted elliptic
genus; equivalently, the Duflot 1970 twisted-character value $4$
for the Mathieu $M_{24}$ class $2A$ restricted to a K3 symplectic
involution). Borcherds lift produces $\Delta_2^{(1)}$ with
$c_2(0) = 4$ and weight $4/2 = 2$. Value
$\kappa_{\mathrm{BKM}}(\Phi_2) = 2$. Matches Family (ii) twined
Borcherds weights in the chapter (eight-form-cover row $N = 2$
primary base); the $N = 2$ multi-valued entry $\{w_N, c_N(0)\}$
= $\{\{2, 4\}, \{4, 8\}\}$ records that the paramodular ring at
level $2$ has two primitive Borcherds bases (Bryan--Oberdieck 2018
Thm.~1), with $(w_N, c_N(0)) = (2, 4)$ being Gritsenko's
$\Delta_2$ and $(4, 8)$ being the non-CHL weight-$4$ companion.

\emph{Verification at $N = 6$ (terminal CHL entry).} The
$g_6$-twined Jacobi form is $\phi_{0, 1}^{(g_6)}(\tau, z) =
(r^{-1} + 2 + r) + O(q)$ with constant $f^{(g_6)}(0, 0) = 2$
(Gaberdiel--Hohenegger--Volpato 2015 Table~1). Borcherds lift
produces $\Delta_1^{(6)}$ with $c_6(0) = 2$ and weight $2/2 = 1$.
Value $\kappa_{\mathrm{BKM}}(\Phi_6) = 1$. Matches.

Cycle~4 closes. The canonical formula $\kappa_{\mathrm{BKM}}
(\Phi_N) = c_N(0)/2$ is valid across the full eight-form spread
at integer, half-integer, and quarter-integer weights, with the
cover assignment $\mathrm{Sp}_4 \to \mathrm{Mp}_4 \to
\widetilde{\mathrm{Mp}}_4$ controlling the weight stratification.

\subsection*{Cycle 5. $\kappa_{\mathrm{BCOV}} = \chi(X)/24$ for
  Class B: first-principles derivation via the BCOV holomorphic
  anomaly}

\paragraph{Attack.} For Class B CY$_3$ targets (quintic, $\C^3$,
resolved conifold, local $\P^2$), the Borcherds-weight identity
$\kappa_{\mathrm{BKM}}$ is undefined because no K3-fibration exists.
The memory prescribes the \emph{replacement} invariant
$\kappa_{\mathrm{BCOV}} = \chi(X)/24$, where $\chi(X) =
\chi_{\mathrm{top}}(X)$ is the topological Euler characteristic.
From what first principle does this replacement follow? The
reflex should be to derive $\chi(X)/24$ from the BCOV holomorphic
anomaly, not to stipulate it as a definition.

\paragraph{Heal.} The invariant $\kappa_{\mathrm{BCOV}} = \chi(X)/24$
is the \emph{genus-$1$ topological-string free energy of BCOV}, which
is forced by the BCOV holomorphic anomaly equation to be
moduli-independent on the BTT moduli space $M_X$, hence a
topological invariant of $X$ alone, hence a function of $\chi(X)$
alone, hence of the form $\chi(X)/c$ for some universal constant
$c$. The constant $c = 24$ is fixed by three independent routes:

(i) \emph{One-loop Quillen determinant}. The BCOV genus-$1$ free
energy on compact CY$_3$ is
$F_1(X) = -\tfrac{1}{2} \log \det(\bar\partial^\dagger
\bar\partial|_X)$, the logarithm of the Quillen norm of the
Hodge-filtered $\bar\partial$-Laplacian (BCOV 1994
\emph{Comm.\ Math.\ Phys.}~165 (1994) 311, \S\S 5.1--5.2). The
$\zeta$-regularised Quillen norm is computed via Grothendieck--
Riemann--Roch applied to the Hodge bundle: at $h^{1, 0} = 0$ the
zero-mode contribution vanishes and the Quillen norm reduces to
a Laplacian determinant on the anti-holomorphic $(0, \bullet)$
forms, whose anomaly coefficient is the Grothendieck--Riemann--Roch
$\int_X \mathrm{ch}(T_X)\,\mathrm{Td}(T_X) = \chi(X)$ by the
CY condition $c_1(T_X) = 0$ collapsing the Todd genus to
$\mathrm{Td}(T_X) = 1 + \tfrac{c_2}{12} + \cdots$. Only the
$(0, 1)$-component survives the $\bar\partial$-compatibility, and
the universal normalisation $1/24$ comes from the
Bott--Chern one-loop $\zeta(0) = -1/2$ regularisation of
$\sum_{n = 1}^{\infty} n = -1/12$, halved by the Quillen norm
convention. The resulting universal one-loop anomaly is
$\alpha_{\mathrm{BCOV}} = (\chi(X)/24) \cdot [\Omega_X]^{0, 1}
\in H^1(X, \mathcal{O}_X)$, the Kapranov--Costello--Li
\emph{arXiv:1606.00365} Prop.~5.2 coefficient.

(ii) \emph{Constant-map contribution}. The BCOV genus-$g$
constant-map free energy on a compact CY$_3$ is
$F_g^{\mathrm{const}}(X) = \chi(X) \cdot c_g$ with
$c_g = (-1)^{g - 1} B_{2g}\, B_{2g - 2}/(2g \cdot (2g - 2) \cdot
(2g - 2)!)$ (Marino--Moore 1998 \emph{Nucl.\ Phys.}~B543,
592; Klemm--Pandharipande 2007 Cor.~1 at $d = 3$; Faber--
Pandharipande 2000 $\lambda_g$-conjecture). At $g = 1$,
$c_1 = 1/24$ by the MacMahon constant-map derivation: the
generating function of constant maps $\sum_{n \geq 0} p(n)\,q^n =
\prod_{n \geq 1} (1 - q^n)^{-1}$ evaluated at genus $1$ via the
tau-function $\log(\eta(q)) = -\tfrac{1}{24} \log q + \cdots$
gives the $\eta$-function tracking factor $1/24$.

(iii) \emph{Gopakumar--Vafa degeneration}. The genus-$1$
constant-map contribution to the Gopakumar--Vafa refined free
energy is $F_1 = \sum_{d \geq 0} n^1_d\, q^d$ with $n^1_0 =
\chi(X)/24$ the leading constant-degree coefficient
(Gopakumar--Vafa 1998 \emph{arXiv:hep-th/9812127}). At the
quintic $X_5$ with $\chi(X_5) = -200$, this gives
$F_1^{\mathrm{const}}(X_5) = -200/24 = -25/3$, matching the
BCOV one-loop evaluation (Candelas--de la Ossa--Green--Parkes
1991 \emph{Nucl.\ Phys.}~B359, 21 at genus $1$).

The three routes (Quillen determinant, constant-map derivation,
Gopakumar--Vafa refinement) converge on $F_1^{\mathrm{BCOV}}(X) =
\chi(X)/24$. This fixes $\kappa_{\mathrm{BCOV}} := \chi(X)/24$ as
the universal Class-B Borcherds-replacement invariant. Examples:

\emph{Quintic $X_5 \subset \P^4$.} $\chi(X_5) = 2(h^{1, 1} -
h^{2, 1}) = 2(1 - 101) = -200$. $\kappa_{\mathrm{BCOV}}(X_5) =
-200/24 = -25/3$.

\emph{$\C^3$ (non-compact toric).} $\chi_{\mathrm{top}}(\C^3) =
1$ in the compactly-supported reading; the one-loop BCOV
coefficient vanishes because $\mathrm{At}(T_{\C^3}) = 0$ by the
Atiyah-retraction to the $T$-fixed origin. In the open-string
BCOV \cite{Costello-Li 1606.00365} reading, the boundary-link
$\beta\gamma$ Casimir energy is $-1/12$; multiplied by the
Euler characteristic gives $\kappa_{\mathrm{BCOV}}(\C^3) = 1/24$
in a certain boundary-regularised reading, or $0$ in the
compactly-supported-interior reading. The invariant is
construction-dependent at non-compact $X$; the two readings
coexist.

\emph{Resolved conifold $X_{\mathrm{con}} = \mathrm{Tot}(\mathcal{O}
(-1)^{\oplus 2} \to \P^1)$.} $\chi_{\mathrm{top}}(\P^1) = 2$
through the retraction $X_{\mathrm{con}} \simeq \P^1$, giving
$\kappa_{\mathrm{BCOV}}(X_{\mathrm{con}}) = 2/24 = 1/12$. The
Costello--Li one-loop BV curving on the boundary link
$T^{1, 1} \simeq S^2 \times S^3$ produces
$\kappa_{\mathrm{ch, BV}}(X_{\mathrm{con}}) = -1$ by the
$\zeta$-regularised $\beta\gamma$ Casimir at $c = -2$
(Polyakov critical-dimension; Kausch \texttt{hep-th/9510149}). The
two invariants differ: BCOV $F_1$ sees the compact $\P^1$,
Costello--Li one-loop curving sees the non-compact boundary.

\emph{Local $\P^2 = \mathrm{Tot}(\omega_{\P^2})$.}
$\chi_{\mathrm{top}}(\P^2) = 3$, giving
$\kappa_{\mathrm{BCOV}}(\mathrm{loc}\,\P^2) = 3/24 = 1/8$. The
$\Z/3$-McKay-quiver reading gives
$\kappa_{\mathrm{ch}}(\mathrm{loc}\,\P^2) = 3/2$
(\texttt{thm:local-p2-shadow}), which is different (by a factor
of $12$) from the BCOV reading. The two invariants
disambiguate: Hodge-supertrace at $d = 3$ is Serre-silent, but the
McKay-quiver substitute on non-compact toric CY$_3$ gives the
shadow value; the BCOV invariant is separately
$\chi/24$ on compact $\P^2 \subset \mathrm{loc}\,\P^2$.

Cycle~5 closes.

\subsection*{Assembled output: five-row $\kappa$-stratification
  at $d = 1, \ldots, 5$ with multi-path verification}

The five cycles produce the following canonical inscription,
extending the chapter by one synoptic summary table. The
verification paths for each $(d, \text{family})$ are annotated.

\begin{center}\small
\begin{tabular}{llccc}
\toprule
$(d, X)$ & $h^{0, \bullet}(X)$ &
  $\kappa_{\mathrm{ch}} = \Xi(X)$ & $\chi(\mathcal{O}_X)$ &
  Verification \\
\midrule
$(1, E)$ & $(1, 1)$ & $0$ & $0$ &
  Serre; $\chi(\mathcal{O}_E) = 1 - g = 0$ \\
$(2, \mathrm{K3})$ & $(1, 0, 1)$ & $2$ & $2$ &
  Griffiths--Harris Ch.~1; Hodge diamond \\
$(2, A_{\mathrm{ab}})$ & $(1, 2, 1)$ & $0$ & $0$ &
  K\"unneth $\Xi(E \times E') = 0 \cdot 0$ \\
$(2, B_{\mathrm{bi}})$ & $(1, 1, 0)$ & $0$ & $0$ &
  $h^{0, 2} = 0$; pairwise cancellation \\
$(3, X_5)$ & $(1, 0, 0, 1)$ & $0$ & $0$ &
  Serre; $\chi(\mathcal{O}_{X_5}) = 0$ odd-$d$ \\
$(3, \mathrm{K3} \times E)$ & $(1, 1, 1, 1)$ & $0$ & $0$ &
  K\"unneth $\Xi(\mathrm{K3}) \cdot \Xi(E)$ \\
$(3, E^3)$ & $(1, 3, 3, 1)$ & $0$ & $0$ &
  K\"unneth $\Xi(E)^3$ \\
$(3, \mathrm{loc}\,\P^2)$ & (open) & $3/2$ & n/a &
  $\tfrac{1}{2}\chi_{\mathrm{top}}(\P^2) = 3/2$;
  $\Z/3$-McKay \\
$(3, X_{\mathrm{con}})$ & (open) & $1$ & n/a &
  Kronecker quiver; $\chi_{\mathrm{top}}(\P^1)/2 = 1$ \\
$(4, X_6)$ & $(1, 0, 0, 0, 1)$ & $2$ & $2$ &
  KP07 Tab.~1; $\chi_{\mathrm{top}} = 2610$ \\
$(4, X_8)$ & $(1, 0, 149, 0, 1)$ & $151$ & $151$ &
  HKTY95 App.~B \\
$(4, \mathrm{K3}^{[2]})$ & $(1, 0, 1, 0, 1)$ & $3$ & $3$ &
  Beauville 1983 + G\"ottsche 1990 \\
$(4, F(Y))$ & $(1, 0, 1, 0, 1)$ & $3$ & $3$ &
  Beauville--Donagi 1985 \\
$(5, X_7)$ & $(1, 0, 0, 0, 0, 1)$ & $0$ & $0$ &
  Serre; strict CY$_5$ \\
$(5, \mathrm{K3}\times X_5)$ & $(1, 0, 1, 1, 0, 1)$ & $0$ & $0$ &
  K\"unneth $\Xi(\mathrm{K3}) \cdot \Xi(X_5) = 2 \cdot 0$ \\
$(5, \text{generic CY}_5)$ & $(1, 0, 0, 0, 0, 1)$ & $0$ & $0$ &
  Serre universal \\
\bottomrule
\end{tabular}
\end{center}

\begin{center}\small
\begin{tabular}{ccc|cccc}
\toprule
$N$ & Cover & $c_N(0)$ & $w_N = c_N(0)/2$ &
  Host CY$_3$ & K3-symp.\ order & Borcherds-status \\
\midrule
$1$ & $\mathrm{Sp}_4(\Z)$ & $10$ & $5$ & $\mathrm{K3} \times E$ &
  $1$ & PROVED \\
$2$ & $\mathrm{Sp}_4(\Z)$ & $4$ & $2$ & $(\mathrm{K3}\times E)/\Z_2$ &
  $2$ & PROVED \\
$3$ & $\mathrm{Sp}_4(\Z)$ & $2$ & $1$ & $(\mathrm{K3}\times E)/\Z_3$ &
  $3$ & PROVED \\
$4$ & $\mathrm{Sp}_4(\Z)$ & $2$ & $1$ & $(\mathrm{K3}\times E)/\Z_4$ &
  $4$ & PROVED \\
$5$ & $\mathrm{Mp}_4(\Z)$ & $1$ & $1/2$ & $(S_5\times E_5)/(\iota_S
  \times \iota_E)$ & $5$ (Borcea-Voisin) & CONJECTURAL (CY$_3$-level) \\
$6$ & $\mathrm{Sp}_4(\Z)$ & $2$ & $1$ & $(\mathrm{K3}\times E)/\Z_6$ &
  $6$ & PROVED \\
$7$ & $\widetilde{\mathrm{Mp}}_4(\Z)$ & $1/2$ & $1/4$ &
  $(\mathrm{K3}\times E)/\Z_7$ (Nikulin-singular) &
  $7$ & OBSTRUCTED (gerbe correction) \\
$8$ & $\mathrm{Sp}_4(\Z)$ & $0$ & $0$ & $\mathrm{Kum}^3(A)/\Z_8$
  (MTW 2018 fourfold) & $8$ & DEGENERATE (abelian limit) \\
\bottomrule
\end{tabular}
\end{center}

The canonical formula $\kappa_{\mathrm{BKM}}(\Phi_N) = c_N(0)/2$ is
universal across the eight-form spread at the full
$\mathrm{Sp}_4$-$\mathrm{Mp}_4$-$\widetilde{\mathrm{Mp}}_4$ cover
stratification; the weight reads off the Borcherds-input principal-
part constant at the cusp via Borcherds 1998 Thm.~13.3 and its
metaplectic refinements. For Class B CY$_3$ (quintic, $\C^3$,
resolved conifold, local $\P^2$), the Borcherds-weight invariant
is undefined, and the replacement is
$\kappa_{\mathrm{BCOV}} = \chi(X)/24$ derived from the BCOV
one-loop Quillen determinant, the genus-$1$ constant-map
contribution, and the Gopakumar--Vafa leading coefficient
$n^1_0 = \chi(X)/24$.

\subsection*{Numerical decomposition refutation: additive BKM split
  fails at all $N \geq 2$}

A tempting (but false) decomposition writes
$\kappa_{\mathrm{BKM}}(\Phi_N) = \kappa_{\mathrm{ch}}(\mathcal{A}_{X_N})
+ \chi(\mathcal{O}_{\mathrm{fiber}_N})$. The $N = 1$ numerical
coincidence reads $\kappa_{\mathrm{BKM}}(\Phi_1) = 5$ while
$\kappa_{\mathrm{ch}}(\mathcal{A}_{\mathrm{K3}\times E}) +
\chi(\mathcal{O}_E) = 0 + 0 = 0 \neq 5$, so even at $N = 1$ the
decomposition fails on the Hodge supertrace reading. The alternative
decomposition with the K3 fibre $\chi(\mathcal{O}_{\mathrm{K3}}) =
2$ plus the Heisenberg-Mukai $\kappa_{\mathrm{ch}}^{\mathrm{Heis}}
(\mathrm{K3}) = 3$ gives $2 + 3 = 5$ at $N = 1$, a numerical
coincidence of two distinct Stage-2 constructions on two distinct
fibrations. At $N \geq 2$, let $X_N := (\mathrm{K3} \times E)/(\Z/
N\Z)$ and $E_N$ its elliptic fibre. The decomposition predicts:
\[
  \kappa_{\mathrm{BKM}}(\Phi_N) \stackrel{?}{=}
  \kappa_{\mathrm{ch}}(\mathcal{A}_{X_N}) +
  \chi(\mathcal{O}_{E_N}).
\]
At $N = 2$: $X_2$ has $\chi_{\mathrm{top}}(X_2) = 0$ and
$\kappa_{\mathrm{ch}}(\mathcal{A}_{X_2}) = 0$ by Hodge supertrace,
while $\chi(\mathcal{O}_{E_2}) = 0$; the right side is $0 + 0 = 0$,
but $\kappa_{\mathrm{BKM}}(\Phi_2) = 2$. Mismatch. At $N = 3$:
right $= 0 + 0 = 0$, left $= 1$. Mismatch. At $N = 4$: right $= 0$,
left $= 1$. Mismatch. At $N = 6$: right $= 0$, left $= 1$. Mismatch
at every $N \in \{2, 3, 4, 6\}$; the decomposition fails uniformly
and is not a structural identity.

The universal formula is $\kappa_{\mathrm{BKM}}(\Phi_N) = c_N(0)/2$
with $c_N(0)$ the Borcherds-input principal-part constant at the
cusp, independent of any decomposition into fibre Euler
characteristics. The coincidence at $N = 1$ on the
Heisenberg-Mukai reading is the \emph{only} place the decomposition
appears to hold, and that is a coincidence of construction
(Heisenberg-Mukai Stage-2 specialisation on $E$ plus K3
structure-sheaf Euler) with the Borcherds weight ($= 5$).

\subsection*{Manin closure}

Five cycles, five closures. The Stage-1 Hodge supertrace
$\kappa_{\mathrm{ch}}(\mathcal{A}_X) = \sum_q (-1)^q\, h^{0, q}
(X)$ is the canonical degree-two shadow coefficient of
$\Phi^{\mathrm{FA}}_d(\mathcal{C})$ on compact CY$_d$, converging
through HKR, Lurie $(\infty, 2)$-trace, and Caldararu Mukai
conventions; it vanishes at odd $d$ by Serre and records
$2, 3, n + 1$ on the three Beauville--Bogomolov strata of even
$d$. The Stage-2 Borcherds-weight identity $\kappa_{\mathrm{BKM}}
(\Phi_N) = c_N(0)/2$ is universal across the eight-form
Gritsenko--Cl\'ery spread with $\mathrm{Sp}_4$-$\mathrm{Mp}_4$-
$\widetilde{\mathrm{Mp}}_4$ cover assignment, reducing to
$(5, 2, 1, 1, 1/2, 1, 1/4, 0)$ and degenerating at $N = 8$. The
Class B replacement $\kappa_{\mathrm{BCOV}} = \chi(X)/24$ is the
BCOV genus-$1$ constant-map coefficient, derived from the
Quillen determinant, MacMahon constant-map, and Gopakumar--Vafa
leading coefficient. The three pillars (Hodge supertrace,
Borcherds weight, BCOV $F_1$) are three distinct construction-
level invariants of a single CY datum; the additive decomposition
$\kappa_{\mathrm{BKM}} = \kappa_{\mathrm{ch}} +
\chi(\mathcal{O}_{\mathrm{fiber}})$ is a numerical coincidence at
$N = 1$ and fails at every $N \geq 2$.

The universal trace identity
$K \circ \Phi = \kappa_{\mathrm{ch}}$ (Vol I $\leftrightarrow$
Vol III bridge) contains these five cycles as its Class-A
evaluations.

\emph{--- End of Opus 12 attack--heal cycles.}
